\chapter{绪论}

\section{研究背景：深度学习对数据的依赖与数据选择的核心地位}

过去十余年间，深度学习在计算机视觉、自然语言处理和多模态理解等领域取得了一系列里程碑式的突破。2012年，AlexNet在ImageNet大规模视觉识别挑战赛\cite{deng2009imagenet}中的优异表现标志着深度学习时代的开启，此后卷积神经网络（Convolutional Neural Network, CNN）迅速成为图像分类、目标检测和语义分割等视觉任务的主流方法\cite{he2016deep, simonyan2015very}。在自然语言处理领域，以GPT系列\cite{radford2019language}和Transformer\cite{vaswani2017attention}为代表的大规模预训练语言模型（Large Language Model, LLM）证明了通过海量文本数据训练的模型可以涌现出强大的语言理解与生成能力。在生成模型方面，扩散模型（Diffusion Model）在图像生成任务上展现了前所未有的质量与多样性。这些突破虽然涉及不同的模型架构和训练范式，但都建立在一个共同的基础之上——大规模高质量数据。

在大模型时代，数据的地位进一步凸显。研究者们逐渐认识到，LLM的许多能力涌现并非仅源于模型规模的增长或架构的改进，而是在很大程度上取决于训练数据的规模、质量与多样性\cite{hoffmann2022training, kaplan2020scaling}。高质量数据的获取与筛选已成为前沿AI研究与工业实践中的核心竞争力\cite{albalak2024survey}。Ilya Sutskever曾发出"数据即将耗尽"的警示——无论这一判断��终是否准确，它深刻反映了当前研究界对数据资源的高度依赖与焦虑。这一趋势也暗示着，单纯追求"更多数据"的粗放式路径正在接近其极限，"更聪明地使用数据"正变得愈发迫切。

然而，数据规模的增长在带来模型能力提升的同时，也引发了一系列现实挑战。首先，计算成本随数据量的增长呈超线性增长\cite{hoffmann2022training}——训练一个大规模模型所需的计算资源已达到普通研究机构难以承受的水平。其次，高质量标注数据日益稀缺且获取成本高昂，尤其在医疗、法律等专业领域，专家标注的成本使得全量标注变得不切实际。此外，大量收集的数据中不可避免地存在冗余、噪声和分布不均等问题，这些低质量数据不仅浪费计算资源，还可能损害模型的泛化性能\cite{sorscher2022beyond}。上述挑战共同指向一个核心认识：\textbf{并非所有数据都同等有价值}。如何从大量候选数据中识别和利用最相关的数据，成为提升深度学习效率与可靠性的关键问题。

数据选择（Data Selection）正是应对这一挑战的核心方法论\cite{albalak2024survey}。广义而言，数据选择是指在深度学习的不同阶段，根据特定目标从候选数据中识别、筛选或评估最相关样本的过程。它不仅涵盖训练阶段常见的核心子集选择（Coreset Selection）\cite{mirzasoleiman2020coresets, killamsetty2021grad}——即从全量训练数据中选取具有代表性的子集用于高效训练，还包括数据混合配比（Data Mixture）——即确定不同来源数据的最优混合比例，数据质量过滤——即剔除低质量或有害样本\cite{chen2023alpagasus, zhou2023lima}，以及推理阶段对输入数据有效性的判断——即判定模型是否应对某一输入做出预测。值得指出的是，数据选择与数据增强（Data Augmentation）���数据集蒸馏（Dataset Distillation）\cite{wang2018dataset, cazenavette2022dataset}有着本质区别：数据选择是从现有数据中"选"，而后两者是通过变换或优化"造"新数据。本文聚焦于数据选择这一方向。

数据选择并非仅局限于深度学习流程中的某一特定环节，而是以不同形式贯穿训练、微调与推理等多个阶段。在\textbf{训练阶段}，核心任务是从全量数据中选出最具信息量的子集，在有限计算资源下最大化模型性能\cite{mirzasoleiman2020coresets, paul2021deep}。在\textbf{微调阶段}，面对来源多样、质量参差的候选数据（如人工标注数据、模型生成数据、开源数据集等），需要选择最匹配下游任务的数据子集或配比，在有限微调预算下高效适配预训练模型\cite{zhou2023lima, chen2023alpagasus}。在\textbf{推理阶段}，模型部署到真实环境后，输入数据的来源不再可控，需要判断每一个输入样本是否处于模型的"有效数据域"内——即模型是否有能力对该输入做出可靠预测\cite{hendrycks2016msp, yang2021generalized_survey}。数据选择在这三个阶段以不同的形式出现、面临不同的约束条件，但都服务于同一个目标：让深度学习系统更高效、更可靠地运行。

进一步地，数据选择的最优策略并非孤立决定的，而是与多种关键因素存在深层交互。\textbf{计算资源约束}影响最优训练数据子集的组成——在不同计算预算下，对模型最有价值的数据可能截然不同\cite{yin2024compute}。\textbf{模型容量约束}与数据冗余之间存在耦合——当模型同时进行压缩（如剪枝\cite{han2015learning}）时，数据的最优选择策略需要随模型状态的变化而调整。\textbf{推理时分布偏移}则要求在推理阶段对输入数据的有效性做出在线判断——训练时的数据选择界定了模型的有效域边界，推理时需要检测输入是否越出这一边界\cite{yang2021generalized_survey}。然而，现有的数据选择研究大多将问题孤立处理：训练数据选择方法不考虑计算预算的影响\cite{killamsetty2021grad, paul2021deep}，数据选择与模型压缩被视为独立的流水线步骤\cite{evci2020rigging}，推理时的分布外检测与训练时的数据选择之间缺乏联系\cite{sun2021react, liu2020energy}。本文旨在填补这一空白，系统研究数据选择与上述关键因素的交互。


\section{研究问题}

基于上述分析，本文从两个互补的视角出发，围绕数据选择提出三个核心研究问题。\textbf{阶段视角}（主线A）关注数据选择在深度学习不同阶段扮演的角色，\textbf{交互变量视角}（主线B）关注数据选择与关键因素之间的交互如何影响选择策略。表~\ref{tab:research-questions}展示了三个研究问题在这两个视角下的定位。

% [表1-1: 三个研究问题的双视角定位]
\begin{table}[htb]
  \centering
  \caption{三个研究问题的双视角定位}
  \label{tab:research-questions}
  \begin{tabular}{cll}
    \toprule
    研究问题 & 主线A：阶段视角 & 主线B：交互变量视角 \\
    \midrule
    子问题一 & 训练/微调阶段的数据选择 & 数据选择 $\times$ 计算资源约束 \\
    子问题二 & 训练阶段的数据与模型联合选择 & 数据选择 $\times$ 模型容量约束 \\
    子问题三 & 推理阶段的数据相关性判断 & 数据选择 $\times$ 推理时分布偏移 \\
    \bottomrule
  \end{tabular}
\end{table}

三个子问题呈递进关系：从阶段视角看，从训练侧（子问题一、二）延伸到推理侧（子问题三）；从交互视角看，交互维度从计算资源到模型容量再到推理时分布偏移，逐步扩展。

\subsection{子问题一：计算预算约束下的训练数据选择}

第一个研究问题聚焦于训练与微调阶段最基础的数据选择问题：\textbf{给定有限的计算预算，如何从全量���练数据中选出使模型性能最优的子集？}

从阶段视角看，这是数据选择在训练阶段的直接体现。从交互视角看，这一问题的核心在于数据选择与计算资源约束之间的交互。已有研究表明，不同计算预算下最优数据子集的组成截然不同\cite{sorscher2022beyond, yin2024compute}：当计算预算有限时，模型仅能学习数据中的简单模式，此时简单、干净的样本更有价值；当计算预算充裕时，模型有能力从多样化甚至含噪声的数据中提取更丰富的信息。然而，现有的数据选择方法普遍忽视计算预算这一关键变量——它们隐含地假设最优数据子集独立于训练时间\cite{killamsetty2021grad, paul2021deep, toneva2018empirical}，导致没有任何单一算法能在不同计算预算下一致领先\cite{xia2024rethinking}。这一问题的核心挑战在��：如何将计算预算作为数据选择框架中的显式变量，使选择策略能够自适应地随预算变化。

\subsection{子问题二：数据选择与模型剪枝的协同}

第二个研究问题同属训练阶段，但将"选择"的维度从数据扩展到模型：\textbf{当训练数据选择与模型剪枝同步进行时，如何利用二者的交互实现协同优化？}

从阶段视角看，模型剪枝（Weight Pruning）是训练阶段的另一种效率优化手段，它通过移除冗余参数来压缩模型\cite{han2015learning, lecun1989optimal}。从交互视角看，数据冗余与参数冗余之间存在双向耦合：一方面，冗余样本（如噪声数据）迫使模型为拟合这些样本而保留不必要的权重，从而阻碍有效剪枝；另一方面，冗余参数扩大了模型的有效表示空间，增加了核心子集选择的搜索难度，导致选出的子集质量下降。然而，现有工作将数据选择和模型剪枝视为独立的流水线步骤——先选数据再训练完整模型，或先训练完整模型再剪枝——未利用也未分析二者之间的耦合关系\cite{evci2020rigging, killamsetty2021grad}。这一问题的核心挑战在于：如何在训练过程中同时进行数据选择与模型剪枝，利用它们的交互建立协同效应。

\subsection{子问题三：推理时的数据有效性判断}

第三个研究问题将数据选择的视角从训练阶段延伸至推理阶段：\textbf{模型部署后，如何判断输入样本是否处于训练数据选择所界定的"有效数据域"内？}

从阶段视角看，模型训练完成并部署后，输入数据的来源不再可控。当输入样本偏离训练数据分布时，模型可能产生过度自信的错误预测\cite{hendrycks2016msp}，这在安全关键应用（如自动驾驶\cite{filos2020can, janai2020computer}、医学诊断\cite{pooch2020can}）中尤为危险。分布外检测（Out-of-Distribution Detection, OOD Detection）可被理解为推理时的数据"选择"——选择（接受）属于训练分布内的输入，拒绝偏离训练分布的输入\cite{yang2021generalized_survey}。从交互视角看，训练阶段的数据选择界定了模型的有效域边界；当推理时输入分布发生偏移，需要判断输入是否仍属于该有效域。这是数据选择与推理时分布偏移（Inference-time Distribution Shift）的交互。这一问题的核心挑战在于：现有的后处理OOD检测方法主要利用全局平均池化后的特征或最终输出层信息\cite{liu2020energy, hendrycks2016msp, sun2021react, sun2022dice}，丢失了池化前通道内的空间激活分布信息，限制了检测能力。


\section{本文贡献概览}

针对上述三个研究问题，本文围绕数据选择这一核心方法论，分别提出一个框架或方法，共同构成对数据选择问题从训练到推理的系统性研究。本文的主要贡献如下：

\textbf{贡献一：计算预算感知的数据选择框架CADS（第三章）。}针对子问题一，本文提出CADS（Computation-Aware Data Selection）框架，首次将计算预算作为双层优化中的显式约束纳入数据选择\cite{zhou2022probabilistic}。在技术层面，CADS通过概率化重参数化和REINFORCE策略梯度实现无Hessian的外层梯度估计，并进一步提出基于惩罚松弛的单层目标转化，将计算开销降低至标准双层优化的3至20倍以下。% [VERIFY]
CADS包含两种变体：CADS-E实现样本级选择，CADS-S实现来源级选择，在视觉分类和语言模型微调任务上，相较基线方法最高提升14.42\%。% [VERIFY]
该工作发表于NeurIPS 2025。

\textbf{贡献二：数据选择与权重剪枝的协同优化框架SWaST（第四章）。}针对子问题二，本文提出SWaST（Simultaneous Weight and Sample Tailoring）框架，通过交替执行权重剪枝与核心子集选择建立协同效应。SWaST首次透明地分析了数据冗余与参数冗余的双向耦合关系，发现并定义了"双重损���"（Critical Double-Loss）问题——即关键参数与其支撑样本被同时移除导致的不可逆退化。为解决这一问题，SWaST引入了状态保护机制，通过知识蒸馏约束防止协同优化过程中的灾难性遗忘。在四个基准数据集和两种网络架构上，SWaST实现了最高17.83\%的准确率提升，同时削减10\%至90\%的浮点运算量（FLOPs）。% [VERIFY]
该工作发表于AAAI 2026。

\textbf{贡献三：基于神经激活先验的分布外检测方法NAP（第五章）。}针对子问题三，本文提出NAP（Neural Activation Prior），发现了一个被现有方法忽视的OOD检测信号：在全连接层前的特征图中，分布内（In-Distribution, ID）样本在各通道内倾向于激活少量神经元且产生较大激活值，而分布外样本的激活则更为分散。基于这一神经激活先验，NAP设计了即插即用的评分函数，无需额外训练，并可与现有OOD检测方法\cite{liu2020energy, hendrycks2016msp, sun2021react, sun2022dice, djurisic2022extremely, sun2022knn}通过加权几何均值进行正交互补。在CIFAR\cite{krizhevsky2009cifar}和ImageNet\cite{deng2009imagenet}基准上，NAP与多种基线方法组合后，FPR95最高降低66.03\%。% [VERIFY]
NAP同时支持CNN和Vision Transformer\cite{dosovitskiy2020image}架构。该工作目前处于投稿中。


\section{论文结构}

图~\ref{fig:research-map}展示了本文三项工作在双主线框架下的位置关系。横轴为阶段递进，从训练侧（第三、四章）延伸到推理侧（第五章）；纵轴为交互维度，从计算资源约束（第三章）到模型容量约束（第四章）再到推理时分布偏移（第五章）。主线A（阶段视角）回答"在深度学习的哪个阶段做数据选择"，主线B（交互变量视角）回答"数据选择与什么关键因素交互"。两条主线相互补充：主线A提供直觉清晰的叙事脉络，主线B揭示工作之间的深层统一性。

% [图1-1: 研究地图——三项工作在双主线框架下的位置。需用户手动绘制。]
\begin{figure}[htb]
  \centering
  % \includegraphics[width=0.8\textwidth]{figures/research-map.pdf}
  \fbox{\parbox{0.8\textwidth}{\centering\vspace{3cm}
    \textbf{占位图：研究地图}\\[0.5em]
    横轴：阶段递进（训练$\to$推理）\\
    纵轴：交互维度（计算资源$\to$模型容量$\to$分布偏移）\\
    三个方框分别对应CADS / SWaST / NAP
  \vspace{3cm}}}
  \caption{研究地图——三项工作在双主线框架下的位置关系}
  \label{fig:research-map}
\end{figure}

本文其余部分的组织如下：

\textbf{第二章\quad 相关工作}，从四个方面综述与本文研究相关的已有工作——训练数据选择方法、模型冗余与效率优化、数据选择中的优化方法、分布外检测，揭示现有研究忽视数据选择与关键因素交互的共性不足。

\textbf{第三章\quad 计算预算感知的训练数据选择}，研究数据选择与计算资源约束的交互。本章从一个探索性实验出发，揭示不同计算预算下最优数据子集的组成截然不同，进而提出将计算预算作为显式变量的数据选择框架CADS。

\textbf{第四章\quad 数据选择与权重剪枝的协同优化}，研究数据选择与模型容量约束的交互。本章首先通过多项式插值实验透明地分析数据冗余与参数冗余的双向耦合，然后提出交替优化框架SWaST，同时削减两类冗余并建立协同效应。

\textbf{第五章\quad 推理阶段的数据相关性判断}，研究数据选择与推理时分布偏移的交互。本章提出神经激活先验（NAP），利用全局池化前通道内的激活分布模式区分分布内与分布外样本，为推理时的数据有效性判断提供新的检测信号。

\textbf{第六章\quad 总结与展望}，从阶段覆盖和交互维度两个视角统一回顾三项工作的贡献，讨论各自的局限性，并展望数据选择领域的未来研究方向。
